\section{Methodology}
\label{sec:methodology}

\subsection{DEM simulation of hollow cylinder tests}
\label{subsec:dem_hca}

DEM simulations were performed in \textsc{PFC3D} \citep{Itasca2021} to reproduce undrained cyclic torsional shear in an HCA. Unlike periodic-cell element tests, the present HCA model uses rigid boundaries (inner and outer cylinders, top and bottom planes) and six torsional blades to apply shear while preserving laboratory-relevant boundary conditions.

The specimen geometry (Fig.~\ref{fig:specimen_generation}) follows a laboratory-scale HCA configuration with inner diameter \SI{6}{cm}, outer diameter \SI{10}{cm}, and specimen height \SI{10}{cm} \citep{Vargas2020}. A rolling-resistance contact model \citep{IwashitaOda1998,WensrichKatterfeld2012} is adopted to capture non-spherical grain effects, with a particle size distribution in the range of \SIrange{1.0}{3.0}{mm}, which is a 10$\times$ scaled-up analogue of the grain size distribution of Toyoura sand. The DEM parameters are summarized in Table~\ref{tab:dem_parameters}. Particles are generated by gravity deposition from the upper part of the apparatus. Gravity is then removed to ensure a uniform stress state, and six vertically arranged torsional blades are inserted into the specimen (Fig.~\ref{fig:specimen_generation_b}). Compared with the conventional approach of constraining boundary particles to transmit torque \citep[e.g.,][]{Li2014,Liu2021}, these independent wall-based blades avoid interference with the radial expansion or contraction of the cylinders, thereby preserving accurate radial stress measurement.

\begin{table}[t]
  \centering
  \caption{Parameters in DEM simulation.}
  \label{tab:dem_parameters}
  \begin{tabular}{@{}lr@{}}
    \toprule
    Description & Value \\
    \midrule
    Number of particles, $N_p$                          & 53{,}764 \\
    Particle density, $\rho_s$ (\si{kg/m^3})            & 2650 \\
    Young's modulus, $E$ (\si{GPa})                     & 0.3 \\
    Normal-to-shear stiffness ratio, $\kappa$           & 2.0 \\
    Interparticle friction, $\mu_{pp}$                  & 0.0\textsuperscript{a}\,/\,0.50\textsuperscript{b} \\
    Particle--wall friction, $\mu_{pw}$                 & 0.0 \\
    Normal critical damping ratio, $\beta_n$            & 0.7 \\
    Shear critical damping ratio, $\beta_s$             & 0.5 \\
    Rolling friction coefficient, $\mu_r$               & 0.1 \\
    \bottomrule
    \multicolumn{2}{@{}l}{\footnotesize \textsuperscript{a}Isotropic consolidation stage.} \\
    \multicolumn{2}{@{}l}{\footnotesize \textsuperscript{b}Anisotropic consolidation and cyclic shear stages.}
  \end{tabular}
\end{table}

\begin{figure}[t]
  \centering
  \begin{subfigure}[t]{\linewidth}
    \centering
    \safeincludegraphics[width=0.5\linewidth]{specimen_generation_pouring.jpg}
    \caption{Particle generation modeling air-pluviation.}
    \label{fig:specimen_generation_a}
  \end{subfigure}
  \vspace{6pt}
  \begin{subfigure}[t]{\linewidth}
    \centering
    \safeincludegraphics[width=0.5\linewidth]{specimen_generation_blade_insertion.jpg}
    \caption{Insertion of torsional blades.}
    \label{fig:specimen_generation_b}
  \end{subfigure}
  \caption{Initial specimen generation process in the DEM-based hollow cylinder apparatus model.}
  \label{fig:specimen_generation}
\end{figure}

The stress and strain evaluation expressions for the HCA geometry are summarized in Appendix~\ref{sec:appendix_hca}.

\subsection{Specimen preparation and anisotropic consolidation protocol}
\label{subsec:preparation}

Specimens are prepared using an isotropic-consolidation to anisotropic-consolidation (IC--AC) protocol. First, the assembly is compressed under zero interparticle friction to a target void ratio, which is determined through the iterative calibration procedure described in Section~\ref{subsec:calibration_validation}: approximately 0.736 for the dense state ($D_r \approx 90\%$) and 0.742 for the medium-dense state ($D_r \approx 75\%$). The interparticle friction coefficient is then restored to 0.5, and the specimen is consolidated from an isotropic state (\Kzero{} = 1.0, $\pmean=\SI{10}{kPa}$) to the target mean effective stress ($\pmean=\SI{100}{kPa}$) along a linear anisotropic stress path.

The cyclic shear simulation program consists of two sets. In the primary set, three representative \Kzero{} states---0.5 (compression), 1.0 (isotropic), and 2.0 (extension)---are examined at a dense state ($D_r \approx 90\%$) across multiple CSR levels to characterize how liquefaction resistance varies with the direction and magnitude of stress anisotropy, and to trace the associated energy dissipation, stiffness degradation, and fabric evolution. In the expanded set, a medium-dense state ($D_r \approx 75\%$) is introduced alongside the dense state, and two additional \Kzero{} values (0.67 and 1.5) are included to refine the coverage of anisotropy levels, yielding ten cases ($2 \times 5$) at a single CSR = 0.300. This expanded set allows examination of whether the \Kzero{}-dependent liquefaction resistance trend observed at $D_r \approx 90\%$ persists at a lower density, and provides a broader parametric basis for the subsequent microscopic analysis.

During the AC stage, $\pmean$ is increased from \SI{10}{kPa} to \SI{100}{kPa} through ten equally spaced target stress states along a linear anisotropic stress path. The resulting stress paths and void-ratio evolution for the three primary cases are shown in Fig.~\ref{fig:ac_path}. The markers in Fig.~\ref{fig:ac_path_a} correspond to the prescribed target states; the solid lines between them represent the actual stress evolution, which exhibits minor fluctuations that arise because different \Kzero{} states require different amounts of axial strain to reach the same $\pmean$ increment---compression-state specimens (\Kzero{} = 0.5) undergo larger axial strain per step, causing the axial stress to lag slightly behind the radial stress during each increment. These fluctuations vanish at each target state, confirming that the servo accurately satisfies the prescribed stress conditions. The final void ratio differs only slightly among the three states (on the order of $10^{-3}$), which allows the comparison to focus on fabric-related effects rather than void-ratio differences.

\begin{figure}[t]
  \centering
  \begin{subfigure}[t]{\linewidth}
    \centering
    \safeincludegraphics[width=0.75\linewidth]{anisotropic_consolidation_stress_path.png}
    \caption{$\pmean$--$q$ stress path.}
    \label{fig:ac_path_a}
  \end{subfigure}
  \vspace{6pt}
  \begin{subfigure}[t]{\linewidth}
    \centering
    \safeincludegraphics[width=0.75\linewidth]{anisotropic_consolidation_void_ratio.png}
    \caption{$\pmean$--$e$ evolution.}
    \label{fig:ac_path_b}
  \end{subfigure}
  \caption{Stress and void-ratio evolution during anisotropic consolidation for the three representative \Kzero{} states.}
  \label{fig:ac_path}
\end{figure}

\begin{figure}[t]
  \centering
  \begin{subfigure}[t]{\linewidth}
    \centering
    \safeincludegraphics[width=0.5\linewidth]{specimen_k0_05_after_ac.jpg}
    \caption{\Kzero{} = 0.5}
    \label{fig:ac_specimens_a}
  \end{subfigure}
  \vspace{6pt}
  \begin{subfigure}[t]{\linewidth}
    \centering
    \safeincludegraphics[width=0.5\linewidth]{specimen_k0_10_after_ac.jpg}
    \caption{\Kzero{} = 1.0}
    \label{fig:ac_specimens_b}
  \end{subfigure}
  \vspace{6pt}
  \begin{subfigure}[t]{\linewidth}
    \centering
    \safeincludegraphics[width=0.5\linewidth]{specimen_k0_20_after_ac.jpg}
    \caption{\Kzero{} = 2.0}
    \label{fig:ac_specimens_c}
  \end{subfigure}
  \caption{Cross-sectional views of DEM HCA specimens after anisotropic consolidation to $\pmean=\SI{100}{kPa}$, colored by particle radius with contact-force vectors superimposed.}
  \label{fig:ac_specimens}
\end{figure}

\subsection{Undrained cyclic torsional shear loading}
\label{subsec:cyclic_loading}

Undrained conditions are represented using the standard constant-volume approximation in DEM, without explicit fluid--particle coupling \citep{YimsiriSoga2010}. The key challenge in simulating undrained HCA tests lies in satisfying multiple boundary conditions simultaneously. The HCA specimen has four kinematic degrees of freedom---inner radius rate ($dr/dt$), outer radius rate ($dR/dt$), height rate ($dH/dt$), and rotation angle rate ($d\theta/dt$)---which must be determined at each timestep to satisfy four conditions: a prescribed radial stress difference, a prescribed axial stress condition, a target shear stress, and constant specimen volume. Earlier DEM-based HCA models adopted simplified servo strategies: \citet{Han2024jgs} solved only the radial and torsional conditions while keeping the specimen height fixed, thereby bypassing the axial degree of freedom entirely; \citet{Ma2024} introduced an axial servo but derived each servo coefficient from the contact stiffness at the corresponding wall independently, without reflecting the cross-coupling among the wall movements. In this study, a fully coupled analytical servo is derived that formulates the boundary conditions directly in terms of effective pressure differences---corresponding to the physically independent control channels of a laboratory HCA---and incorporates the constant-volume constraint to yield servo coefficients that simultaneously satisfy all four conditions (radial pressure difference, axial pressure difference, constant volume, and cyclic torsional shear) with the radial--axial interaction accounted for exactly. The servo controls two pressure quantities measured from particle--wall contact forces: the radial pressure difference $p_{dif,r}^{\prime}=p_o^{\prime}-p_i^{\prime}$ and the axial pressure difference $p_{dif,z}^{\prime}=p_z^{\prime}-\sigma_r^{\prime}$ (see Table~\ref{tab:hca_equations} for definitions). 
\paragraph{Radial condition.}
In a laboratory HCA, the outer and inner cell pressures are applied independently through separate pressure lines; the radial condition reproduces this capability by controlling the pressure difference between the outer and inner cylinders. The radial pressure difference and its servo error are defined as:
\begin{equation}
p_{dif,r}^{\prime} = p_o^{\prime} - p_i^{\prime}, \qquad e_r = p_{dif,r}^{\prime} - p_{dif,r}^{tar},
\label{eq:cond1_radial}
\end{equation}
where $p_o^{\prime}$ and $p_i^{\prime}$ are the outer and inner radial effective pressures (Table~\ref{tab:hca_equations}) and $p_{dif,r}^{tar}$ is the target radial pressure difference. The inner cylinder radius rate is adjusted at each timestep to eliminate the radial error:
\begin{equation}
\frac{dr}{dt} = \frac{e_r}{\Delta t}\,S_{cr},
\label{eq:servo_radial}
\end{equation}
where $\Delta t$ is the simulation timestep and $S_{cr}$ is the radial servo coefficient---a gain factor, derived from the aggregate contact stiffness at the cylindrical walls, that converts the pressure error into the wall velocity needed to eliminate it. Its explicit expression is given in the coupled analytical servo formulation below (\cref{eq:Scr_constantH_main}).

\paragraph{Axial condition.}
Laboratory HCA tests can operate under either constant total axial stress or constant specimen height. To reproduce the former, the servo must control the difference between the vertical and radial effective stresses, analogous to the radial pressure difference in \cref{eq:cond1_radial}. The axial pressure difference and its servo error are defined as:
\begin{equation}
p_{dif,z}^{\prime} = p_z^{\prime} - \sigma_r^{\prime}, \qquad e_z = p_{dif,z}^{\prime} - p_{dif,z}^{tar},
\label{eq:cond2_axial}
\end{equation}
where $p_z^{\prime}$ is the vertical effective pressure (Table~\ref{tab:hca_equations}), $\sigma_r^{\prime}$ is the average radial effective stress, and $p_{dif,z}^{tar}$ is the target axial pressure difference. The specimen height rate is adjusted accordingly:
\begin{equation}
\frac{dH}{dt} = \frac{e_z}{\Delta t}\,S_{cz},
\label{eq:servo_axial}
\end{equation}
where $S_{cz}$ is the axial servo coefficient. Two operating modes are implemented:
(i)~\emph{Stress-control mode}: $p_{dif,z}^{tar}$ is prescribed by the loading protocol, and the coupled servo (\cref{eq:servo_dr,eq:servo_dH}) adjusts both $dr/dt$ and $dH/dt$ simultaneously. This mode is used during anisotropic consolidation and monotonic shear validation (\cref{subsec:calibration_validation}), where the axial stress condition must be actively maintained to match the experimental protocol of \citet{Nakata1998}.
(ii)~\emph{Displacement-control mode}: the specimen height is held constant ($dH/dt=0$), and $p_{dif,z}^{\prime}$ evolves freely. This mode is adopted for all cyclic shear simulations to match standard HCA cyclic test procedures \citep{Ishihara1985}.

\paragraph{Torsional condition.}
The torsional servo is independent of the radial--axial system and follows the same formulation adopted in earlier DEM-based HCA studies \citep{Han2024jgs,Ma2024}. The rotation angle rate is controlled to achieve the target shear stress:
\begin{equation}
\frac{d\theta}{dt} = \frac{T - T^{tar}}{\Delta t}\,S_{cs},
\label{eq:cond3_torsion}
\end{equation}
where $T$ is the current torque, $T^{tar}$ is the target torque, $\Delta t$ is the simulation timestep, and $S_{cs}$ is the torsional servo coefficient. The coefficient $S_{cs}$ is determined from the contact-stiffness-weighted rotational stiffness $I_{rot}$ about the specimen axis (Fig.~\ref{fig:servo_geometry}):
\begin{equation}
I_{rot} = \sum r_d^2\,k_n\cos^2\!\theta, \qquad S_{cs} = \frac{1}{I_{rot}},
\label{eq:irot_scs}
\end{equation}
where $r_d$ is the distance from the rotation axis to the contact point, $k_n$ is the normal contact stiffness, and $\theta$ is the angle between the contact normal and the horizontal plane.

\begin{figure}[t]
  \centering
  \safeincludegraphics[width=0.5\linewidth]{servo_torsional_contact_geometry.png}
  \caption{Contact geometry used to evaluate the rotational stiffness contribution for torsional servo control. $r_d$: distance to rotation axis; $\theta$: angle between contact normal and horizontal.}
  \label{fig:servo_geometry}
\end{figure}

\paragraph{Volume constraint.}
The undrained (constant-volume) condition requires (with the outward-positive sign convention $dr>0$, $dR>0$ denoting radius increases):
\begin{equation}
2\pi H\!\left(R\frac{dR}{dt} - r\frac{dr}{dt}\right) + \pi\!\left(R^2-r^2\right)\frac{dH}{dt} = 0.
\label{eq:const_volume_constraint}
\end{equation}

\paragraph{Analytical servo.}
Although \cref{eq:servo_radial,eq:servo_axial} are written in uncoupled form for clarity, the radial and axial conditions are in fact coupled through two mechanisms: (i) the constant-volume constraint (\cref{eq:const_volume_constraint}) links $dR/dt$ to $dr/dt$ and $dH/dt$, so that an inner-cylinder wall movement alters the outer radius; and (ii) the axial pressure difference $p_{dif,z}^{\prime}=p_z^{\prime}-\sigma_r^{\prime}$ depends on $\sigma_r^{\prime}=({p_o^{\prime}R+p_i^{\prime}r})/({R+r})$ (Table~\ref{tab:hca_equations}), which is a weighted average of the inner and outer pressures affected by both cylinder wall movements. An analytical solution is obtained by reducing the four unknowns to two through two successive eliminations:
\begin{enumerate}
\item \emph{Torsion.} The torsional servo (\cref{eq:cond3_torsion}) does not enter the kinematic constraints that govern $dr/dt$, $dR/dt$, and $dH/dt$. It is therefore solved independently at each timestep, leaving three unknowns ($dr/dt$, $dR/dt$, $dH/dt$) and three conditions.
\item \emph{Volume constraint.} \Cref{eq:const_volume_constraint} expresses $dR/dt$ as a function of $dr/dt$ and $dH/dt$, reducing the system to two unknowns ($dr/dt$, $dH/dt$) with two target conditions ($e_r=0$ and $e_z=0$ in \cref{eq:cond1_radial,eq:cond2_axial}).
\end{enumerate}

Let $K_{r,i}$ and $K_{r,o}$ denote the aggregate normal contact stiffness at the inner and outer cylinders, $A_i=2\pi rH$ and $A_o=2\pi RH$ the corresponding wall areas, and $K_z=\frac{1}{2}\sum_{c\in\mathrm{top/bottom}} k_n$ the effective axial contact stiffness with area $A_z=2\pi(R^2-r^2)$. Substituting \cref{eq:const_volume_constraint} into the pressure increments yields the coupled stiffness relation (derived in Appendix~\ref{sec:appendix_servo}):
\begin{equation}
\begin{pmatrix} \delta p_{dif,r}^{\prime} \\[4pt] \delta p_{dif,z}^{\prime} \end{pmatrix}
=
\underbrace{\begin{pmatrix} \hat{K}_{11} & \hat{K}_{12} \\[6pt] \hat{K}_{21} & \hat{K}_{22} \end{pmatrix}}_{\hat{\mathbf{K}}}
\begin{pmatrix} dr \\[4pt] dH \end{pmatrix},
\label{eq:coupled_stiffness}
\end{equation}
where
\begin{alignat}{2}
\hat{K}_{11} &= -\frac{K_{r,i}}{A_i}-\frac{r}{R}\,\frac{K_{r,o}}{A_o}, &\qquad
\hat{K}_{12} &= \frac{R^2-r^2}{2RH}\,\frac{K_{r,o}}{A_o}, \label{eq:Khat_main_row1}\\[4pt]
\hat{K}_{21} &= -\frac{K_{r,i}-\frac{r}{R}\,K_{r,o}}{A_i+A_o}, &\qquad
\hat{K}_{22} &= -\frac{K_z}{A_z}-\frac{(R^2-r^2)\,K_{r,o}}{2RH\,(A_i+A_o)}. \label{eq:Khat_main_row2}
\end{alignat}
The entries carry dimensions of Pa/m (see Appendix~\ref{sec:appendix_servo} for the full derivation). Three entries are negative, reflecting the outward-positive convention, while $\hat{K}_{12}>0$ arises from the volume-constraint-induced cross-coupling. The determinant of $\hat{\mathbf{K}}$ is
\begin{equation}
\det(\hat{\mathbf{K}}) = \hat{K}_{11}\,\hat{K}_{22} - \hat{K}_{12}\,\hat{K}_{21}.
\label{eq:det_K_main}
\end{equation}
Setting $(\delta p_{dif,r}^{\prime},\,\delta p_{dif,z}^{\prime})=(-e_r,\,-e_z)$ and inverting \cref{eq:coupled_stiffness} yields the coupled servo equations that replace the uncoupled forms of \cref{eq:servo_radial,eq:servo_axial}:
\begin{align}
\frac{dr}{dt} &= \frac{1}{\det(\hat{\mathbf{K}})}\!\left(-\hat{K}_{22}\,\frac{e_r}{\Delta t} + \hat{K}_{12}\,\frac{e_z}{\Delta t}\right), \label{eq:servo_dr}\\[6pt]
\frac{dH}{dt} &= \frac{1}{\det(\hat{\mathbf{K}})}\!\left(\hat{K}_{21}\,\frac{e_r}{\Delta t} - \hat{K}_{11}\,\frac{e_z}{\Delta t}\right), \label{eq:servo_dH}
\end{align}
The remaining kinematic variable $dR/dt$ is then obtained by substituting $dr/dt$ and $dH/dt$ into \cref{eq:const_volume_constraint}. A relaxation factor less than unity is applied to all servo coefficients to ensure numerical stability. In the present study, equal inner and outer cell pressures are adopted ($p_{dif,r}^{tar}=0$), and the constant-height condition ($dH/dt=0$) is used for all cyclic shear simulations, so that only the radial condition remains active and the servo coefficient reduces to
\begin{equation}
S_{cr} = -\frac{1}{\hat{K}_{11}} = \frac{1}{\dfrac{K_{r,i}}{A_i}+\dfrac{r}{R}\,\dfrac{K_{r,o}}{A_o}}.
\label{eq:Scr_constantH_main}
\end{equation}

\paragraph{Cyclic loading protocol.}
Specimens are subjected to sinusoidal torsional shear under the constant-height condition ($dH/dt=0$). The cyclic stress ratio (CSR) is defined as the maximum applied torsional shear stress normalized by the initial mean effective stress. In the primary analyses, liquefaction is identified using a stress-based criterion, $\ru = 0.95$, where $\ru$ is the excess pore pressure ratio. A strain-based criterion (single-amplitude shear strain $\gamma_{SA}=2.5\%$) is adopted in the calibration/validation stage (\cref{subsec:calibration_validation}) to match the criterion used by \citet{Ishihara1985}. A comparison of the two criteria across all simulated cases confirms that they yield nearly identical $N_L$ values and the same resistance ordering among \Kzero{} states, so the stress-based criterion is used for all subsequent analyses.

The excess pore water pressure is inferred from the reduction in radial effective stress:
\begin{equation}
u = \sigma_{r0}^{\prime} - \sigma_r^{\prime},
\label{eq:epwp_definition}
\end{equation}
where $u$ is the excess pore water pressure, $\sigma_r^{\prime}$ is the current radial effective stress, and $\sigma_{r0}^{\prime}$ is its initial value. The ratio $\ru$ is then defined as $\ru = u/\sigma_{r0}^{\prime}$.

\subsection{Parameter calibration and validation}
\label{subsec:calibration_validation}

The micro-parameters listed in Table~\ref{tab:dem_parameters} were calibrated against the undrained monotonic torsional shear data of \citet{Nakata1998} for dense Toyoura sand ($D_r \approx 90\%$, $\pmean_0=\SI{100}{kPa}$, \Kzero{} = 1.0) \citep[see also][]{Ma2024}. Because the micro-parameters collectively influence shear stiffness, dilatancy, and liquefaction resistance, the calibration targeted these behaviors in sequence. Starting from a base parameter set at a given void ratio, the contact Young's modulus $E$ and rolling friction coefficient $\mu_r$ were adjusted while keeping the normal-to-shear stiffness ratio $\kappa$ constant, aiming to reproduce the shear stiffness observed in the monotonic shear tests. The effective stress path in $\pmean$--$\tau_{z\theta}$ space was then examined to assess dilatancy. If the simulated dilatancy was insufficient---i.e., the mean effective stress at phase transformation was lower than the experimental value---the specimen was further compacted to a lower void ratio, which increased both dilatancy and shear stiffness; $E$ was then reduced accordingly to restore the target shear stiffness. This iterative process between void ratio, contact modulus, and rolling friction was repeated until both the shear stiffness and dilatancy simultaneously matched the monotonic experimental data. During the monotonic calibration, the axial pressure difference was maintained at zero ($p_{dif,z}^{tar}=0$ in \cref{eq:cond2_axial}) to ensure $\sigma_z^{\prime}=\sigma_r^{\prime}$ during shearing.

The calibrated parameter set was then validated against the cyclic liquefaction resistance data of \citet{Ishihara1985} for the same sand under identical conditions. Specimens were subjected to sinusoidal torsional shear at five CSR levels (0.20, 0.25, 0.30, 0.35, and 0.40), with liquefaction defined by single-amplitude shear strain $\gamma_{SA}=2.5\%$.

\begin{figure*}[t]
  \centering
  \safeincludegraphics[width=\textwidth]{validate_combined.png}
  \caption{Calibration and validation of the DEM-HCA model for dense Toyoura sand ($D_r \approx 90\%$, $\pmean_0=\SI{100}{kPa}$, \Kzero{} = 1.0): (a) effective stress path and (b) stress--strain relationship from undrained monotonic shear, calibrated against \citet{Nakata1998}; (c) cyclic stress ratio versus number of cycles to liquefaction, validated against \citet{Ishihara1985}.}
  \label{fig:validation_combined}
\end{figure*}

Fig.~\ref{fig:validation_combined} summarizes the calibration and validation results. The effective stress path (Fig.~\ref{fig:validation_combined}a) captures the initial contraction followed by dilation towards the critical state line, matching the experimental trend. The stress--strain relationship (Fig.~\ref{fig:validation_combined}b) reproduces the shear stiffness and strain-hardening behavior. These two aspects confirm a successful calibration against the monotonic data. The liquefaction resistance curve (Fig.~\ref{fig:validation_combined}c) shows good agreement with the cyclic data across the full range of CSR, validating the calibrated parameters under cyclic loading conditions without further adjustment. The simultaneous agreement across stiffness, dilatancy, and cyclic liquefaction resistance confirms the reliability of both the micro-parameters (Table~\ref{tab:dem_parameters}) and the combined servo mechanism described in \cref{subsec:cyclic_loading}.

Having calibrated and validated the model under isotropic consolidation (\Kzero{} = 1.0), specimens with different \Kzero{} values obtained from the IC--AC protocol were subjected to undrained cyclic torsional shear. Unlike the monotonic calibration where $p_{dif,z}^{\prime}$ was actively controlled, constant height ($dH/dt=0$) is adopted for all cyclic shear simulations to match typical HCA test protocols.

\subsection{Microscopic descriptors}
\label{subsec:micro_descriptors}

While the macroscopic response characterizes the overall liquefaction behavior, a key advantage of DEM is its ability to access particle-scale information that is not directly measurable in laboratory tests. To exploit this advantage, three microscopic descriptors are introduced to probe the contact network from complementary perspectives: the mechanical coordination number \Zm{}, the signed fabric anisotropy indicator $\alpha$, and the contact density $\rho_c$.

The mechanical coordination number \citep{Thornton2000,Oda1977} is defined as
\begin{equation}
Z_m = \frac{2N_c - N_1}{N_p - N_1 - N_0},
\label{eq:zm_definition_method}
\end{equation}
where $N_c$ is the total number of interparticle contacts, $N_p$ is the total number of particles, and $N_1$ and $N_0$ are the numbers of particles with one and zero contacts, respectively. Floaters are excluded to focus on the load-bearing contact network. $Z_m$ serves as a scalar measure of packing compactness and contact network connectivity.

The directionality of the contact network is characterized through the fabric tensor \citep{Oda1982fabric,Oda1982stat}:
\begin{equation}
\Phi_{ij} = \frac{1}{N_c}\sum n_i\,n_j,
\label{eq:fabric_tensor}
\end{equation}
where $n_i$ and $n_j$ are the components of the contact normal vector. Because the HCA specimen is axisymmetric about the vertical axis $z$, it is natural to express $\Phi_{ij}$ in the cylindrical coordinate system $(r,\theta,z)$; the axisymmetric stress condition ($\sigma_r = \sigma_\theta$) ensures that $z$ is a principal direction of the fabric tensor, so the axial diagonal component $\Phi_{zz}$ is a principal value. The signed scalar indicator $\alpha$ \citep{YimsiriSoga2010,Otsubo2022} is then defined from $\varphi_z \equiv \Phi_{zz}$:
\begin{equation}
\alpha = \frac{5(3\varphi_z - 1)}{5\varphi_z + 1}.
\label{eq:alpha_definition_method}
\end{equation}
Positive $\alpha$ indicates that contact normals concentrate toward the vertical axis (compression-state anisotropy, \Kzero{} $< 1.0$), whereas negative $\alpha$ indicates concentration toward the horizontal plane (extension-state anisotropy, \Kzero{} $> 1.0$).

The contact density \citep{Han2023} describes the average number of contacts per unit solid angle on the unit sphere of contact normal orientations:
\begin{equation}
\rho_c(\theta_z,\,\phi_{\mathrm{cir}}) = \frac{2\,N(\theta_z,\,\phi_{\mathrm{cir}})}{N_p\,\Delta\Omega},
\label{eq:contact_density}
\end{equation}
where $\theta_z$ and $\phi_{\mathrm{cir}}$ are the polar and azimuthal angles in the spherical coordinate system, $N(\theta_z,\,\phi_{\mathrm{cir}})$ is the number of contacts whose normals fall within the angular bin of widths $\Delta\theta_z$ and $\Delta\phi_{\mathrm{cir}}$, and $\Delta\Omega = (\cos\phi_{\mathrm{cir}} - \cos(\phi_{\mathrm{cir}}+\Delta\phi_{\mathrm{cir}}))\,\Delta\theta_z$ is the solid angle subtended by that bin. Noting that the fabric tensor (\cref{eq:fabric_tensor}) is the second-order moment of the contact normal distribution, $\rho_c$ retains the full directional information---including higher-order features that the fabric tensor projects out---and thus provides a more complete picture of the evolving contact network.

Together, these three descriptors probe the contact network at increasing levels of directional detail: \Zm{} quantifies the overall connectivity (scalar), $\alpha$ captures the dominant anisotropy direction (signed scalar), and $\rho_c$ resolves the full directional distribution (field on the unit sphere). They are used alongside force-chain visualizations to connect particle-scale fabric evolution to the macroscopic liquefaction resistance trends reported in \cref{sec:results}.
